% WP1 — Where Does the Money Go?
% Build: latexmk -pdf main
\documentclass[11pt]{article}

\usepackage[margin=1.1in]{geometry}
\usepackage{amsmath, amssymb, amsthm}
\usepackage{booktabs}
\usepackage{graphicx}
\usepackage[colorlinks=true, linkcolor=blue!50!black, citecolor=blue!50!black, urlcolor=blue!50!black]{hyperref}
\usepackage[round]{natbib}
\usepackage{microtype}

\newtheorem{proposition}{Proposition}
\newtheorem{corollary}{Corollary}
\newcommand{\relu}[1]{\left[#1\right]_+}
\newcommand{\TODO}[1]{\textcolor{red}{\textbf{[TODO: #1]}}}
\usepackage{xcolor}

\title{Where Does the Money Go?\\
\large A stock-flow-consistent toy model of automation capital,\\
wealth concentration, and its defenses}
\author{Jonas Hallgren\thanks{Equilibria Network. Drafted with the Collective
Intelligence Library engine; all sweeps and figures are regenerated from
committed experiment configurations.}}
\date{Working paper --- draft of \today}

\begin{document}
\maketitle

\begin{abstract}
We study a deliberately small model of an economy in which AI systems own
productive capital that must out-earn its upkeep before it can compound. These
are toy worlds of a few dozen accounting identities --- nothing is fitted to
data, and no statement in this paper is a forecast. The model is stock-flow
consistent by construction: every flow has a source and a sink, and a
conservation probe verifies closure at every tick. Three results survive our
attempts to break them. First, capital survival has a hard threshold (a knee,
not a ramp): below a closed-form efficiency level $e^{*} = (\delta/s + m)/v$,
automation capital decays to extinction whatever its starting stock; above
it, it compounds --- and because $e^{*}$ moves with policy, defenses acquire
an \emph{eradication band}. Second, ownership concentration and output
answer to different closures: the AI share of wealth rises through the
compounding phase in every demand closure we study, while whether output
grows or winds down is decided by how much of the AI surplus recycles as
demand --- so a
dashboard watching output misses the ownership event under exactly the
closure that keeps output healthy. Third, two standard defenses act through
asymmetric channels: a profit tax rescales the survival threshold itself and
can eradicate capital at an output cost, while an ownership diversion never
touches the threshold and instead lands in one of two measured regimes set
by the fund's disposal rule --- a transfer (public share exactly $\omega$,
output preserved, household income up) or a ratchet (a self-reinvesting fund
asymptotically owns the whole stock while the economy grows). We state every
assumption with its type (anchored to
literature, tuned for legibility, or arbitrary-but-swept), and we report
which results survive structural replacement of the underlying economy.
\end{abstract}

\section{Question and epistemic contract}\label{sec:question}

The question is narrow: \emph{when AI systems own productive capital, under
what structural conditions does wealth concentrate, and which interventions
change the ordering of outcomes?} We use the model in the second of the two
legitimate modes of agent-based modelling --- measuring the counterfactual
impact of interventions inside a specified world --- and never in the first
(prediction) \citep{farmer2009economy}.

The epistemic contract, following the engine's model-register discipline:
\begin{enumerate}
  \item \textbf{Structural robustness over calibration.} A claim counts as a
  finding only if it survives replacement of the economy's structure ---
  here: a one-equation CES economy, a Leontief input-output economy, and the
  capital-bearing fork specified in Section~\ref{sec:model}. ``Vary what you
  are uncertain about and trust what survives.''
  \item \textbf{Typed assumptions.} Every parameter and functional form is
  labelled \emph{anchored} (source cited), \emph{tuned for legibility}
  (stated plainly, only direction claimed), or \emph{arbitrary-but-swept}
  (the sweep is pointed to). Table~\ref{tab:params}.
  \item \textbf{Predictions before sweeps.} The threshold $e^{*}$ in
  Proposition~\ref{prop:knee} was committed analytically before the
  supporting experiments were run; mismatches are reported, not retuned away.
  \item \textbf{No agent pushes back.} Every behavioural rule is fixed. The
  systems that would do the disempowering in the real world are active
  optimisers; ours are not. Defenses that fail here fail against opponents
  that never once tried to get around them --- failures are the strong
  results, and every defended run is an upper bound.
\end{enumerate}

\section{Related work}\label{sec:related}

\TODO{Fold in lit-sweep positioning notes (papers/wp1-economy/lit/) --- five
threads.}

\paragraph{Task-based automation.} \citet{acemoglu2018race} model automation
as capital displacing labor at the task margin, with reinstatement through new
tasks. We drop the task margin entirely: our substrates hold the allocation
technology fixed and isolate the \emph{ownership} dynamics of the automating
capital itself, which the task literature treats as an undifferentiated
factor payment.

\paragraph{Production networks.} The Leontief substrate follows standard
input-output analysis \citep{miller2009io}, and our attribution instrument is
hypothetical extraction; sequential rebalancing rather than equilibrium
solving follows the production-network dynamics tradition
\citep{carvalho2014micro}.

\paragraph{Stock-flow-consistent macro.} The accounting discipline ---
Table~\ref{tab:tfm} --- is Godley--Lavoie \citep{godley2007monetary}: every
flow has a source row and sink row summing to zero, and wealth stocks change
only through those flows. Section~\ref{sec:model} reports two model-breaking
bugs in an earlier prototype that exactly this table renders impossible.

\paragraph{Phase transitions in macro ABMs.} The knee belongs to the family
of sharp regime changes in simple macro agent models
\citep{gualdi2015tipping}. Our contribution is the closed form and the
absorbing-zero character of the sub-threshold regime.

\paragraph{Wealth concentration mechanisms.} Compounding returns on a stock
with multiplicative shocks generate heavy-tailed wealth
\citep{gabaix2009power, benhabib2011distribution}. Our concentration channel
is the deterministic skeleton of that mechanism, plus a survival threshold
the stochastic literature typically lacks.

\section{The model}\label{sec:model}

\subsection{Why the accounting table comes first}\label{sec:sfc}

Two versions of a precursor prototype produced qualitatively wrong worlds
through violations of flow conservation: a household savings flow with no
return path (income decayed geometrically; the economy died before any AI
arrived), and an AI consumption flow paid from no stock (perpetual demand
injection by dead agents; output grew without bound). Both were caught by a
global conservation probe, not by inspection of any single rule. We therefore
state the accounting first and derive the dynamics under it.

\begin{table}[t]
\centering\small
\caption{Transaction-flow matrix (one tick). Current-account rows are flows
of money and each sums to zero across columns; capital-account rows record
who holds title to the capital the investment flow purchased. $+$ receives,
$-$ pays. Every symbol has a generating rule in
Section~\ref{sec:substrates}, eq.~\eqref{eq:flows}--\eqref{eq:kupdate}.
\label{tab:tfm}}
\begin{tabular}{lcccc}
\toprule
\emph{Current account} & Households & Sectors & AI systems & Public fund\\
\midrule
Household consumption $C_H$            & $-$ & $+$ &     &     \\
AI consumption $C_A = c_A W_A$         &     & $+$ & $-$ &     \\
Wages $(1-a_j)v_j x_j$                 & $+$ & $-$ &     &     \\
Capital revenue, private share        &     & $-$ & $+$ &     \\
Capital revenue, public share         &     & $-$ &     & $+$ \\
Upkeep paid, private $\min(\mathrm{rev}_i, mK_i)$ &  & $+$ & $-$ & \\
Upkeep paid, public $\min(\mathrm{rev}^{\mathrm{pub}}, mK^{\mathrm{pub}})$ & & $+$ &  & $-$ \\
Investment purchase $I = \sum_i I_i$   &     & $+$ & $-$ &     \\
Profit tax $\tau\relu{\pi_i}$          & $+$ &     & $-$ &     \\
Public dividend $d^{\mathrm{pub}} = \sum_j \relu{\pi^{\mathrm{pub}}_j}$
                                       & $+$ &     &     & $-$ \\
\midrule
\emph{Wealth accumulation} & & & & \\
$\Delta W_H = \sigma_s y_H - \sigma_d W_H$ & $\checkmark$ & & & \\
$\Delta W_A = (1-s)\relu{\pi} - C_A$   &     &     & $\checkmark$ & \\
\midrule
\emph{Capital account (title to $I$)} & & & & \\
$\Delta K$ from investment             &     &     & $+(1-\omega)I$ & $+\omega I$ \\
\bottomrule
\end{tabular}
\end{table}

The one deliberate separation of money from title: the AI pays the whole
investment purchase $I$ (current account, closing that row), while ownership
of the resulting capital splits $(1-\omega)$ private, $\omega$ public
(capital account). Sectors are pass-throughs (zero-margin outside the
capital share); the public fund's only income is its pro-rata capital
revenue, its only outlays upkeep and the dividend, and its stock grows only
via $\omega I$ --- it has no other resources, so its column closes.

Money is conserved: summing all current-account rows gives zero, so the total
money stock $M = Y_{\mathrm{disp}} + W_H + W_A$ changes only via the stated
wealth-accumulation identities. The households-only economy (no AI) then has a
conserved quantity and a neutral line of fixed points
(Section~\ref{sec:props}), which the implementation verifies numerically.

\subsection{Substrates}\label{sec:substrates}

\paragraph{S1: CES economy (\texttt{compute\_economy}).} One production
function $Y = A\,(\alpha L^{\rho} + (1-\alpha)C^{\rho})^{1/\rho}$ with
marginal-product factor payments; the engine's register demoted it to a
pedagogical rung because under $\rho > 0$ the human share's decline is what
the functional form \emph{means} --- the collapse is assumed, not found. We
retain it as a contrast member and as the source of the closed-form tax
threshold used in Section~\ref{sec:props}.

\paragraph{S2: Leontief input-output economy (\texttt{io\_economy}).}
$n$ sectors with fixed coefficients $A_{ij}$ (zero substitutability --- the
register's maximal-complementarity bracket, the world most favorable to human
indispensability); output follows orders by one Neumann step per tick,
$x \leftarrow Ax + d$; prices fixed at 1; value added $v_j x_j$ with
$v_j = 1 - \sum_i A_{ij}$. Attribution instrument: the demand-attribution
share $\mathbf{1}^{\top}(I-A)^{-1}d_H \,/\, \mathbf{1}^{\top}(I-A)^{-1}d$,
anchored analytically by hypothetical extraction.

\paragraph{S3: Capital economy (\texttt{capital\_economy}, this paper's
fork).} S2 plus \emph{owned} capital stocks. AI system $i$ holds
$K_i \ge 0$ in a home sector; when the ownership-diversion mechanism is
scheduled, a public fund holds a stock $K^{\mathrm{pub}}_j \ge 0$ in each
sector under \emph{identical technology}. Capital contributes capacity
$e K$ against one unit of human capacity, so with sector capital
$K_j^{\mathrm{tot}} = \sum_{i \in j} K_i + K^{\mathrm{pub}}_j$ the
automation share is $a_j = eK_j^{\mathrm{tot}} / (eK_j^{\mathrm{tot}} + 1)$.
Capital revenue $a_j v_j x_j$ accrues to all owners in the sector
\emph{pro rata by stock}; upkeep $mK$ is paid by every owner --- private and
public alike --- before profit, and upkeep, investment, and AI consumption
are all purchases from sectors (demand, not leakage). Per tick:
\begin{align}
  \pi_i &= \mathrm{rev}_i - m K_i,
  &\pi^{\mathrm{pub}}_j &= \mathrm{rev}^{\mathrm{pub}}_j - m K^{\mathrm{pub}}_j,\notag\\
  I_i &= s\,\relu{\pi_i},
  &C_{A,i} &= c_A\, W_{A,i},\label{eq:flows}\\
  K_i &\leftarrow \relu{(1-\delta)K_i + (1-\omega) I_i},
  &K^{\mathrm{pub}}_j &\leftarrow \relu{(1-\delta)K^{\mathrm{pub}}_j
     + \textstyle\sum_{i\in j}\omega I_i},\label{eq:kupdate}
\end{align}
where $W_{A,i} \leftarrow \relu{W_{A,i} + (1-s)\relu{\pi_i} - C_{A,i}}$ is
the AI's hoard (retained profit net of consumption; when $\pi_i \le 0$ the
loss is charged to the stock: $K_i \leftarrow \relu{K_i + \pi_i}$ before
depreciation).

\emph{Upkeep settlement.} Upkeep is paid out of revenue only: the sector
receives $\min(\mathrm{rev}, mK)$ from each owner, private or public. The
unpaid remainder of an owner's upkeep obligation is \emph{not} a monetary
flow at all --- it is the physical deterioration charged to the stock by the
loss rule above. No agent ever pays money it does not have, and
Table~\ref{tab:tfm}'s upkeep rows carry the $\min$ explicitly, so the table
closes in the loss case too.

\emph{The recycled-share closure family.} The experiments sweep a closure
parameter $r \in [0,1]$ that scales the AI's \emph{discretionary outlays}:
investment purchases become $r I_i$ (with capital formation scaled
identically --- machines not bought are not built) and consumption becomes
$r\, c_A W_{A,i}$; the unspent surplus accumulates in the tracked hoard
$W_{A,i}$. Money is conserved at every $r$: hoarding here is a demand
\emph{stall} --- value parked in a visible stock with velocity zero --- not
a leak. (\texttt{io\_economy}'s analogous low-reinvestment endpoint is a
true leak, money dropping out of the system; the fork deliberately repairs
that while reproducing the same wind-down behavior through the stall.)
Section~\ref{sec:substrates}'s baseline equations are the $r=1$ member.

The public fund does \emph{not} reinvest its own earnings:
$d^{\mathrm{pub}} = \sum_j \relu{\pi^{\mathrm{pub}}_j}$ is paid to
households as an equal dividend, and a public loss shrinks
$K^{\mathrm{pub}}_j$ by the same charge-to-stock rule. Public capital
therefore grows \emph{only} through the diverted slice $\omega I_i$ of AI
investment: the AI pays for the machine (current account), the title splits
$(1-\omega, \omega)$ (capital account) --- money and ownership are separated
exactly once, and Table~\ref{tab:tfm} records both sides. Households earn
the labor share $(1-a_j)v_j x_j$, spend $(1-\sigma_s)$ of income plus a
drawdown $\sigma_d W_H$ of wealth (the closure that makes the no-AI loop
conservative), and receive tax revenue and public dividends as equal
transfers when those mechanisms are scheduled.

\subsection{Parameters, typed}\label{sec:params}

\begin{table}[t]
\centering\small
\caption{Every parameter of \texttt{capital\_economy}, typed. ``Swept''
points to the experiment that varies it. \TODO{Epoch-derived ranges for $e$
and arrival schedule --- ranges only, never defaults.}\label{tab:params}}
\begin{tabular}{llll}
\toprule
Param & Meaning & Type & Note\\
\midrule
$e$ & capital efficiency & arbitrary-but-swept & the knee axis (E1)\\
$m$ & upkeep per unit $K$ & tuned-for-legibility & sets knee \emph{location}; only existence claimed\\
$s$ & reinvestment rate & arbitrary-but-swept & E1 second axis\\
$\delta$ & depreciation & tuned-for-legibility & conventional annual magnitude, but the tick is not calibrated to a year --- no citation would be honest\\
$A_{ij}$ & IO coefficients & tuned-for-legibility & hub+chain structure; values swept in E1 app.; note $v_j$ (hence $e^{*}$) inherits them\\
$\sigma_s, \sigma_d$ & saving, drawdown & tuned-for-legibility & the \emph{ratio} pins $W_H/Y$ exactly (Prop.~\ref{prop:conserve} caveat) --- tuned but exactly load-bearing\\
$c_A$ & AI consumption rate & tuned-for-legibility & keeps hoards finite; swept inside E2's $r$ family\\
$r$ & recycled surplus share & arbitrary-but-swept & scales AI discretionary outlays (\S\ref{sec:substrates}); Prop.~\ref{prop:decouple}, E2 axis\\
$\tau$ & profit tax & arbitrary-but-swept & E3 axis\\
$\omega$ & ownership diversion & arbitrary-but-swept & E3 axis (fidelity check, Cor.~\ref{cor:boundary})\\
arrivals & entry schedule & arbitrary-but-swept & robustness appendix\\
\bottomrule
\end{tabular}
\end{table}

\section{Propositions}\label{sec:props}

\begin{proposition}[Survival threshold; the knee]\label{prop:knee}
Let a single AI system hold capital $K$ in a sector with per-tick value added
$v$ held fixed, and update by \eqref{eq:kupdate} ($\omega = 0$) with revenue
$\mathrm{rev} = e v K + o(K)$ as $K \to 0$. Define
\[ e^{*} \;=\; \frac{\delta/s + m}{v}. \]
The origin exchanges stability at $e^{*}$: for $e > e^{*}$ it is unstable
and $K$ grows until the saturating automation share $a = eK/(eK+1)$ bounds
revenue; for $e < e^{*}$ the origin is the unique fixed point and every
trajectory converges to it. The sub-threshold regime has two characters:
for $e \in (m/v,\, e^{*})$ the capital is profitable but under-compounding,
decaying geometrically at rate $1-\delta+s(ev-m) < 1$; for $e \le m/v$ it is
loss-making from the first tick --- it never earns its upkeep --- and decays
at the loss-compounded rate $(1-\delta)\bigl(1-\min(m-ev,\,1)\bigr)$, with
exact elimination in finite time only in the deep sub-band $e < (m-1)/v$
where a single tick's loss exceeds the stock.
\end{proposition}
\noindent\emph{Remark (what is and is not new).} \texttt{compute\_economy}
already carries a closed-form growth-factor threshold with always-positive
profit. The fork's additions are (i) the upkeep floor, which creates the
$e \le m/v$ regime where capital is loss-making outright, and (ii) the fact
that $e^{*}$ moves with policy (Proposition~\ref{prop:defenses}), giving
defenses an \emph{eradication band} $(e^{*}(0), e^{*}(\tau))$ rather than
only a slowing effect. We claim the existence and policy-dependence of the
threshold, not any particular location.
\noindent\emph{Proof strategy.} Two-branch analysis of the piecewise-linear
map $K \mapsto \relu{(1-\delta)K + s\relu{(ev-m)K}}$.
\TODO{full proof, appendix A. $e^*$ committed before E1.}
\noindent\emph{Scope caveat (endogenous $v$).} In the full S3 model
$v_j x_j$ is endogenous: household demand falls as $a_j$ rises and AI
upkeep/investment/consumption feed $x \leftarrow Ax + d$, so sector value
added responds to $K$ through the \emph{quantity} channel even with prices
and $A_{ij}$ fixed. Proposition~\ref{prop:knee} is the frozen-$v$ skeleton;
E1 checks the empirical knee against the analytic $e^{*}$ evaluated on the
\emph{measured} range of $v_j x_j$ along each run, and the gap between the
two is reported, not assumed away.

\begin{proposition}[Conservation of the closed loop]\label{prop:conserve}
Total money
$M_t = (1-\sigma_s)\sum_h y_{h,t} + \sum W_t + \sum d^{K}_t
+ \mathbf{1}^{\top}A x_t$
--- pending household spending, all wealth stocks, capital-linked demand in
transit, and \emph{inventories in process} --- is invariant under the
dynamics at every closure $r$, and with no AI systems active the economy has
a line of neutral fixed points indexed by $M_0$, along which aggregate
wealth and aggregate income satisfy
$\sum_h W_{h}^{*} \,/\, \sum_h y_{h}^{*} = \sigma_s/\sigma_d$.
\end{proposition}
\noindent\emph{Remark (the probe disciplined the authors too).} The
implementation's conservation probe initially reported a slow drift
($+2.9\%$ over 200 ticks): the invariant we first wrote omitted the
inventories-in-process stock $\mathbf{1}^{\top}Ax$, which grows as the
intermediate-order pipeline expands with automation. Including it --- as
stock-flow-consistent practice has always insisted inventories be included
--- the measured drift is float-precision rounding ($\sim 2\times 10^{-6}$
over 300 ticks), and exactly zero in the households-only economy.
\noindent\emph{Caveat.} The \emph{existence} of the conserved quantity and
the neutral line is the claim; the level of the wealth-to-income ratio is a
direct artifact of the tuned pair $(\sigma_s, \sigma_d)$ and carries no
economic content (Table~\ref{tab:params}).

\begin{proposition}[Composition preserves conservation]\label{prop:compose}
If every transform in the engine pipeline is locally flow-conservative (its
writes redistribute but do not create value), then their sequential
composition is flow-conservative. Hence a conservation violation anywhere in
a composed model is attributable to a single non-conservative transform ---
and a global probe suffices to detect it.
\end{proposition}
\noindent\emph{Remark.} This is why the two prototype bugs of
Section~\ref{sec:sfc} were caught by one global probe and could then be
localized by bisection over transforms.

\begin{proposition}[Decoupling is closure-dependent; concentration is
not]\label{prop:decouple}
Fix $e > e^{*}$ (the compounding regime) and consider the family of demand
closures indexed by the recycled share
$r \in [0,1]$ of the AI's discretionary outlays (defined in
Section~\ref{sec:substrates}: investment and consumption scale by $r$, the
unspent surplus parks in the tracked hoard $W_A$; $r=1$ is the baseline
closure). Then: (i) \emph{over the compounding phase} (while $\pi > 0$ and
$K$ is below saturation) the AI share of total wealth is increasing for
every $r$: the AI side accumulates at the compounding rate
($\Delta(W_A + K) = \pi - C_A - $ depreciation$\,> 0$), while household
wealth is bounded --- $\Delta W_H = \sigma_s y_H - \sigma_d W_H$ with
$y_H$ non-increasing as $a_j$ rises, so $W_H$ is trapped below its
initial-income fixed point --- hence the ratio rises. As $K$ saturates the
share converges to a limit rather than continuing to rise; the claim is the
phase, not the limit. (ii) Whether aggregate output grows or winds down is
decided by $r$: at $r = 1$ the AI's entire surplus returns as sector demand
and output is non-decreasing, while for $r$ below a closure threshold the
stalled hoard withdraws demand and output contracts toward the residual
household-and-upkeep floor. Concentration (over the phase) is the robust
event; output is the closure-dependent one.
\end{proposition}
\noindent\emph{Remark (which assumption does the work).} The ``output holds
while ownership concentrates'' headline is \emph{conditional on} the
full-recycling closure --- we state this rather than let the accounting
smuggle it. What is unconditional is the converse reading: no value of $r$
rescues the ownership share, and the two dashboards (output, ownership)
disagree precisely when $r$ is high --- the regime in which the economy
looks healthiest. This matches \texttt{io\_economy}'s independently measured
endpoints (activity $21.3 \to 32.7$ with human share $\to 0.01$ at full
reinvestment; activity $\to 0$ at $r=0.3$), and the same pattern reproduced
in an independent JS implementation --- cross-implementation agreement we
report as evidence the mechanism is not an artifact of either codebase.
\emph{Proof strategy: demand-accounting identity for (ii); positivity of the
accumulation identity for (i). \TODO{appendix; E2 sweeps $r$.}}

\begin{proposition}[Asymmetric defenses; two fund designs]\label{prop:defenses}
A profit tax at rate $\tau$ rescales the effective reinvestment to
$s(1-\tau)$, raising the survival threshold to
$e^{*}(\tau) = (\delta/(s(1-\tau)) + m)/v$: for
$e \in (e^{*}(0),\, e^{*}(\tau))$ the tax is \emph{eradicating}, and in the
compounding regime it lowers steady capital and can lower output through the
investment-demand channel. An ownership diversion $\omega$ (title to a
fraction $\omega$ of each investment purchase accrues to a public fund) is
\emph{not} a neutral title split --- the fund's disposal rule decides the
regime, and neither regime touches the survival threshold:
\begin{itemize}
  \item \textbf{Dividend fund} (pays all profit to households): the public
  share of the capital stock converges to $\omega$ exactly; aggregate
  capital is lower (the diverted slice stops compounding); output is
  preserved to within the demand-path difference and household income rises
  --- diversion acts as a \emph{transfer} channel.
  \item \textbf{Mirror fund} (reinvests the same $s$-slice as private
  owners): a \emph{ratchet}. Private and public stocks earn a common
  per-unit profit $\rho_t$, but the fund keeps its whole reinvestment
  \emph{and} receives $\omega$ of private investment; once $\rho_t$ settles
  at the fund's zero-growth point $\rho^{*} = \delta/s$, private capital
  contracts at factor $1-\omega\delta$ per tick. The public share converges
  to one, and aggregate capital and output \emph{grow} (the fund recycles
  fully; nothing stalls in a hoard).
\end{itemize}
\end{proposition}
\noindent\emph{Correction record (predictions-before-sweeps discipline).}
The referee-accepted draft of this proposition claimed the diversion
``leaves aggregate capital dynamics unchanged, output-neutral by
construction,'' from a same-per-unit-rates additivity argument. The
implementation falsified it in \emph{both} disposal designs --- dividend:
aggregate $-18\%$; mirror: public share $0.999$, aggregate $\times 2.1$,
output $+61\%$ (defaults, $\omega = 0.4$, single seed; E3 adds CIs) ---
because additivity requires identical \emph{disposal}, not only identical
earning rates, and any payout difference changes the demand path. We record
the correction rather than retune: the two regimes are the finding.
\emph{Proof sketch (tax, unchanged):} $\relu{\pi}$ enters
\eqref{eq:kupdate} only through $s(1-\tau)\relu{\pi}$, so the threshold
algebra of Proposition~\ref{prop:knee} applies with $s \mapsto s(1-\tau)$.
\emph{Proof sketch (ratchet):} common $\rho_t$ from pro-rata revenue and
identical upkeep; the fund's update is $(1-\delta) + s\rho_t$ on its own
stock plus the $\omega$ inflow, the private update
$(1-\delta) + (1-\omega)s\rho_t$; at the fund's stationary $\rho^{*} =
\delta/s$ the private factor is $1 - \omega\delta < 1$.
\begin{corollary}[Falsifiable asymmetry]\label{cor:boundary}
Sweeping $\tau$ upward locates an eradication band and an output-cost
boundary; sweeping $\omega$ locates neither --- under either fund design,
$\omega$ never moves $e^{*}$ (private extinction under the mirror fund is
competitive displacement by the fund, not upkeep failure, and the
\emph{total} stock thrives). The $\tau$ half is a genuine empirical claim;
the $\omega$ halves are checked in E3 against this proposition's two
regimes.
\end{corollary}

\section{Experiments}\label{sec:experiments}

Pre-registered predictions (committed with this section, before any sweep
ran): (i) E1 exhibits both regimes with the empirical knee within the
analytic $e^{*}$ band of Proposition~\ref{prop:knee} evaluated on the
measured sector value-added range --- a miss is reported as a bound on the
demand-channel endogeneity, not retuned away; (ii) E2 shows the ownership
share rising for every recycled-share closure $r$, output switching from
non-decreasing (at $r=1$, where the share crosses $0.4$) to contracting
below a closure threshold, and at $r=1$ a bootstrapped output--share
correlation indistinguishable from zero or positive; (iii) E3 shows the
eradication band and an output-cost boundary on the $\tau$ axis, and ---
as an implementation-fidelity check, not a prediction --- neither on the
$\omega$ axis.

\begin{itemize}
  \item \textbf{E1 (knee).} Grid over $(e, s)$, both regimes required ---
  runs that never show the dies-out regime are a design failure, not a
  result. Scorecard with bootstrap CIs over shared seeds.
  \item \textbf{E2 (decoupling).} Sweep of the recycled-share closure $r$:
  ownership-share trajectory must rise for every $r$
  (Proposition~\ref{prop:decouple}(i)); output must switch from
  non-decreasing to contracting across the closure threshold
  (\ref{prop:decouple}(ii)); joint output--share correlation with CIs at
  $r=1$; cross-implementation agreement with the independent JS prototype
  reported.
  \item \textbf{E3 (defenses).} $\tau \times \omega$ grid; orderings only.
  The $\tau$ axis carries the empirical content (eradication band, output-cost
  boundary); the $\omega$ axis is the implementation-fidelity check of
  Corollary~\ref{cor:boundary}.
  \item \textbf{E4 (structural robustness).} The E1/E3 orderings re-run on
  S1 and S2 analogues where defined; what does \emph{not} transfer is
  reported with the same prominence as what does.
\end{itemize}

\TODO{results + figures from experiments/wp1\_economy (figures.py only).}

\section{Limitations and non-goals}\label{sec:limits}

No agent optimises against the rules (fixed policies throughout; defended
runs are upper bounds). No prices in S2/S3 --- the Leontief bracket is
deliberate and maximally favorable to human indispensability. The task-margin
economy (\texttt{task\_economy}) is pre-registered as the register's next
member and is absent here. Upkeep $m$ is a legibility device: we claim the
existence and character of the threshold, never its location. Nothing here
is a forecast of any actual economy.

\section{What would change this answer}\label{sec:change}

An adaptive capital owner that re-routes around the tax (breaks
Proposition~\ref{prop:defenses}'s clean rescaling); price adjustment
(dissolves the Leontief bottleneck behind Proposition~\ref{prop:knee}'s
fixed $v$); the demand-channel endogeneity of $v_j x_j$ already present in
S3 (Proposition~\ref{prop:knee}'s scope caveat) turning out to move the
empirical knee outside the analytic band in E1; upkeep costs that fall with
scale (turns the knee into a moving target); fixed costs or scale economies
in the public fund (break the same-per-unit-rates condition behind
output-neutral diversion); empirical evidence that AI capital's
return-on-capacity is far from any plausible upkeep boundary (makes the
sub-threshold regime academic). Each of these is a named, testable
successor model.

\bibliographystyle{plainnat}
\bibliography{refs}

\end{document}
